\subsection{Three objects with cyclic support at nineteen arrows}

\begin{lemma}[Proved: detecting a three-by-three strong pair in one corner]
\label{CP2:A-strong-pair-three-three-detector}
Let $S$ be a two-object strong category with trivial minimal group
and total minimal row and column counts both three. If
$H_2(BS;\mathbb Q)\ne0$, the unique object with two minimal rows
is also the unique object with two minimal columns. At that object
$a$, the corner inclusion induces a surjection
$H_2(B\operatorname{End}(a);\mathbb Q)\to H_2(BS;\mathbb Q)$.
\end{lemma}
\begin{proof}
The augmentation modules $V,W$ of
Lemma~\ref{CP2:A-high-Rees} are each one-dimensional. The left
module is supported at the object with two rows, and the right
module at the object with two columns. If these objects differ,
orthogonality of their identity idempotents gives
$W\otimes_A V=0$. Otherwise both are supported at $a$ and
their ordinary tensor product over the full corner maps onto
$W\otimes_A V$. The two augmentation sequences used in the
degree-two band--Tor formula are natural for the corner inclusion,
so this is the asserted map on homology.

Explicitly, label the four minimal arrows at $a$ by $e_{ij}$,
$i,j\in\{0,1\}$. The normalized bar cycle
\[
 [e_{01}\mid e_{10}]-[e_{00}\mid e_{10}]
 -[e_{01}\mid e_{00}]+[e_{00}\mid e_{00}]
\]
has zero boundary, since all four displayed products equal
$e_{00}$. Under dimension shifting it represents the tensor of
the column difference and the row difference. These differences
span $W$ and $V$, respectively. Its image therefore spans the
nonzero quotient $W\otimes_A V$, which also verifies the claimed
naturality directly in degree two.
\end{proof}

\begin{lemma}[Proved: binary stalks into a two-row band corner]
\label{CP2:A-binary-two-row-sink}
Let a finite category be a collage of a source block $S$ and a
sink monoid $N$ whose proper minimal ideal is a rectangular band
with two rows. Suppose $H^2(BS;\mathbb Q)=0$ and every connecting
stalk is nonempty and has at most two states. Every rational
degree-two product on the collage is zero.
\end{lemma}
\begin{proof}
A state fixed by all of $N$ makes the corner inclusion nullhomotopic.
Its restriction and the source-block restriction of every second
class are then zero, so the two-block lemma applies.
Assume that no stalk has such a state. If the transformation
image contains no constant, it consists entirely of permutations.
Every minimal-band element is then an idempotent permutation, hence
maps to the identity; $nk\in K$ forces every $n\in N$ to map to
the identity as well, giving fixed states. If the image contains
exactly one constant, every actor fixes its value. Thus both
constants occur; the image may also contain the transposition.
The minimal band maps onto this two-constant ideal. Its constant
value depends only on its row, and the two row values are distinct.
The stalk is consequently the canonical left act $Ne$ for a
minimal idempotent $e$, with its two row labels.

Every map between two such stalks preserves both row labels.
The source block therefore acts identically on this common
two-element set. The bridge is $B(N,Ne)\times BS$.
The left act $Ne$ is a retract of the regular left act $N$;
the latter action category has initial object $1_N$.
Consequently $B(N,Ne)$ is contractible, and the bridge projection
to $BS$ is a homotopy equivalence. The actual pushout identifies
the collage with $BN$ up to homotopy. Every balanced form for a
two-row minimal band has rank at most one, so
Lemma~\ref{A:rank-one-square} kills its rational second squares.
Polarization proves the assertion for all second products.
\end{proof}

\begin{lemma}[Proved: a nine-arrow strong pair with a ternary P-five stalk]
\label{CP2:A-nine-pair-P5-five-stalks}
Let a strong source pair $S$ have nine arrows, trivial minimal
group, total rows three and total columns two. Attach a sink
$P_5$ with stalk sizes two and three. The resulting nineteen-arrow
category cannot have type $\CP^2$.
\end{lemma}
\begin{proof}
The source pair has finite-group type, in fact is contractible,
because it has one minimal column at each object. If a stalk has
a $P_5$-fixed state, the two-block fixed-arrow argument applies.
Assume neither does. A left $P_5$ act on two or three states with
no fixed state has exactly one two-element core orbit, labelled
by the two minimal rows. Indeed every minimal orbit is a quotient
of the two-element minimal left ideal; a singleton would be fixed.
Every band actor is constant with its row value on the core.
If a third state is present, idempotence forces its image under
that actor to be the same row value. Thus on either stalk the
four band actors are the two row constants. Every intertwiner
preserves the two row-labelled core states.

Write $a$ for the two-row object of $S$ and $b$ for its one-row
object. Both have one minimal column. The category consists of
its six minimal arrows, its two identities and one extra arrow.
Unless the extra arrow is $t:a\to b$, all incoming arrows at
$b$ from the other object are minimal and fixed by its minimal
idempotent $e_b$. There are no external entrances. Its units fix
these arrows as well: the only possible extra endomorphism is
at $b$, where $e_b$ is a zero. The one-sided deletion applies.

Suppose the extra arrow is $t:a\to b$. If the stalk $F(b)$ is
binary, every right endomorphism of $b$ acts identically on it,
since it commutes with the two left row constants. Every internal
outgoing arrow from $b$ is also fixed by $e_b$, since $b$ has one
minimal column. The outgoing one-sided deletion applies.

The only remaining arrangement is $|F(a)|=2$, $|F(b)|=3$.
All right endomorphisms of $a$ act identically on $F(a)$.
Let $u:a\to b$ be the unique minimal arrow of that type, and
choose a minimal idempotent $e_a$ at $a$. Since $te_a=u$,
contravariant functoriality gives $F(t)=F(u)$.
The endomorphism $F(e_b)$ factors through the binary set $F(a)$;
it is the retraction onto the two labelled core states of $F(b)$.
Define $E_a$ to be the identity and $E_b=F(e_b)$.

These maps are a natural idempotent endomorphism of the entire
connecting profunctor. At the two corners the assertion follows
from the identities just given. Every minimal arrow $b\to a$
induces the unique core inclusion $F(a)\to F(b)$, and is fixed
after $E_b$. The map $F(u)$ is fixed after precomposition by
$E_b$, because $e_b$ factors through that core inclusion; the
same holds for $F(t)=F(u)$. These are all arrows of $S$.
The maps also commute with $P_5$, whose actors are row constants.
The image cut thus gives an actual categorical retract, with one
fewer arrow. The cut-retract theorem preserves actual projective-plane
type on this image. Its eighteen-arrow budget is impossible.
\end{proof}

\begin{theorem}[Proved: cyclic three-object support through nineteen]
\label{CP2:A-strong-three-nineteen}
A finite category on three objects with cyclic support and at
most nineteen arrows cannot have classifying space of type $\CP^2$.
\end{theorem}
\begin{proof}
An actual object-retract deletion contradicts the complete
eighteen-arrow exclusion. Assume none is available. A wholly
strong three-object category has finite-group type unless both
total row and column counts are at least four. In the latter
case the nineteen-arrow bound forces the pure $4\times4$
categorical band with three identities. Its degree-four homology
is zero by Lemma~\ref{CP2:A-pure-categorical-band}.

There remains a strong pair $S$ and a single corner $N$. By
opposition take $S$ to be the source. Write $q=|S|$, $m=|N|$
and $h$ for the sum of the two nonempty stalk sizes. We may
assume $q+m+h=19$, $q\ge6$, and $h\ge3$: the smaller budget
is already excluded, while $h=2$ is a singleton join.
If both block second cohomology groups vanish, use the two-block
lemma.

First suppose $H^2(BS;\mathbb Q)\ne0$. Then $q\le15$, the
minimal group is trivial, and the possible total row/column
counts are $(3,3)$, $(3,4)$ and $(4,3)$. If $m\le2$ and
each local row count is at most two, exact sink orbits followed
by the categorical binary projective-cover theorem kill the
second products. The only other local row distribution is
$(1,3)$ with total columns three, forcing $q\ge14$. At $q=14$
the pair is pure; its one-row object admits the incoming one-sided
deletion. At $q=15$, the budget forces $m=1,h=3$, so both stalks
are binary and the categorical binary-act lemma applies directly.

If $m\ge3$, then $q\le13$ and the totals are $(3,3)$.
The pure case $q=11$ again admits the one-row-object deletion.
At $q=12$ or $13$, the small-binary sink lemma applies unless
\[
 q=12,\qquad m=3,\qquad h=4,
\]
and the stalk sizes are one and three. By the detector lemma,
the object $a$ with two rows also has two columns; the other
object $b$ has one of each. If the stalk at $a$ is singleton,
the corner inclusion at $a$ is nullhomotopic in the full collage.
Its surjection on second homology makes every global second class
restrict trivially to $S$. The sink corner has order three and
zero second cohomology, so the two-block lemma applies.

Otherwise the stalk at $b$ is singleton. There is just one
internal arrow beyond the nine-element minimal band and two
identities. If it is not an arrow $a\to b$, all incoming
arrows at $b$ are minimal; $e_b$ fixes them and its units do
as well, since either there are no extra units or $e_b$ is a
corner zero. The incoming one-sided deletion applies. If the
extra arrow is $a\to b$, all internal outgoing arrows at $b$
are minimal and fixed by $e_b$, and the external outgoing
stalk is singleton. The outgoing one-sided deletion applies.
This finishes the positive-source-block case.

Finally suppose $H^2(BS;\mathbb Q)=0$ and
$H^2(BN;\mathbb Q)\ne0$. Then $m\ge5$. A singleton stalk
would give the fixed-arrow obstruction, so $h\ge4$ and $q\le10$.
If the total row count is two, the actual source minimal-group
collapse replaces $S$ by its minimal group. Under this budget
the group is trivial unless $q=10$, and that exceptional budget
will be treated below. The resulting ordinary $N$ cone has
$m\le9$ and bounded homology, so the bounded-eleven act theorem
applies. In all other cases with $q\le9$, the minimal group is
trivial and the row/column totals are $(3,2)$, or at $q=10$
possibly $(4,2)$. For $q=8$ the pair is pure and its one-row
object admits the incoming deletion.

At $q=9$ the only new budgets are $(m,h)=(5,5)$ and $(6,4)$.
The first has $N=P_5$ and is the preceding ternary-stalk lemma.
In the second both stalks are binary. A six-element monoid with
nonzero second cohomology has a proper $2\times2$ minimal band,
so Lemma~\ref{CP2:A-binary-two-row-sink} applies. At $q=10$
the budget forces $m=5,h=4$, so $N=P_5$ and both stalks are
binary. Lemma~\ref{CP2:A-P5-binary-source-block} applies
regardless of the minimal group or the internal shape of $S$.
These are all possibilities.
\end{proof}
